\documentclass[10pt,a4paper]{amsart}
\usepackage[margin=19mm]{geometry}
\usepackage{amsmath,amssymb,mathtools}
\usepackage[scr=rsfs,cal=euler]{mathalfa}
\usepackage{xfrac}
\usepackage[unicode,hidelinks,bookmarks=false]{hyperref}
\newcommand{\ie}{i.e.\ }
\newcommand{\cf}{cf.\ }
\newcommand{\resp}{respectively\ }
\newcommand{\Subsection}{\subsection}
\providecommand{\nobreakdash}{\nobreak\hbox{-}}
\setlength{\emergencystretch}{2em}
\multlinegap0pt
\usepackage{amssymb}
\usepackage{xcolor}
\usepackage{float}
\usepackage{bm}
\newtheorem{proposition}{Proposition}
\title{Stability of Einstein 4-manifolds satisfying a chiral curvature condition}
\author{Diego Artacho}
\date{Source: arXiv:2609.15159v1 (CC BY 4.0)}
\begin{document}
\maketitle

\noindent\emph{Source and licence.} Stability of Einstein 4-manifolds satisfying a chiral curvature condition. Version arXiv:2609.15159v1 (CC BY 4.0), distributed by arXiv under the Creative Commons Attribution 4.0 International licence, http://creativecommons.org/licenses/by/4.0/, as recorded in the arXiv licence metadata for this exact version and shown on the abstract page https://arxiv.org/abs/2609.15159v1. Full licence notice: \texttt{licenses/arxiv-2609.15159-CC-BY-4.0.txt}.\par\smallskip

\noindent\textit{Editorial scope.} Counted unit: the article's proposition prop:spin (source0.tex lines 225--232). The article's complete original proof of it is delivered with it (source0.tex lines 234--236, one \texttt{proof} environment of 3 lines). No mathematics is written, repaired or re-derived: the statement and the proof are the article's own lines, in the article's own notation, and no retained line is substituted or re-flowed.  No further source-local context is needed: every cross-reference of every retained line resolves inside the two delivered spans.  Nothing was omitted from any retained span, so the package carries no omission record.  No retained line contains a citation, so the wrapper carries no bibliography and no bibliography entry.  No retained line contains a figure environment, an image-inclusion command or drawing code, so no image is delivered and the package declares no asset dependency.  Editorial formatting only, in addition: an AMS \texttt{amsart} preamble that restates the article's own declarations and macros used by the retained lines, namely the environments \texttt{proposition} on separate counters, together with the standard packages those lines need.  The retained environments print in this document as Proposition 1 in order of appearance, whereas the article prints the same items with the numbering of its own full document.  The complete native source of the article as distributed in the arXiv e-print for this exact version is bundled unchanged as \texttt{sources/210393/source0.tex}.
\begin{proposition}\label{prop:spin}
    Let $(M,g)$ be an oriented Riemannian four-manifold. Then, the following hold: 
    \begin{enumerate}
        \item $TM \oplus \Lambda^+$ is a spin vector bundle, and the locally defined spinor bundles $\Sigma M = \Sigma^+ M \oplus \Sigma^- M$ of $M$ and $\Sigma \Lambda^+$ of $\Lambda^+$ define an honest spinor bundle $\Sigma M \otimes_{\mathbb{C}} \Sigma \Lambda^+$over $M$; 
        \item $\Sigma^+M \otimes_{\mathbb{C}} \Sigma \Lambda^+ \cong \Sigma^+M \otimes_{\mathbb{C}} \Sigma^+M \cong \mathrm{End}_{\mathbb{C}}(\Sigma^+M)$; and 
        \item $\mathrm{Id}_{\Sigma^+M} \in \Gamma(\mathrm{End}(\Sigma^+M))$ is parallel with respect to the connection induced by the Levi-Civita connection. 
    \end{enumerate}
\end{proposition}

\begin{proof}
    $TM$ and $\Lambda^+$ are associated to the representations $\mathrm{Id}_{\mathrm{SO}(4)}$ and $\widehat{\mathrm{pr}}_+$ of $\mathrm{SO}(4)$ respectively. Hence, locally, their spinor bundles are associated to the representations $\Delta_4 = \Delta^+_4 \oplus \Delta^-_4$ and $\Delta_3 \circ \mathrm{pr}_+$ of $\mathrm{Spin}(4)$ respectively, where $\Delta_n$ denotes the spin representation of $\mathrm{Spin}(n)$. However, note that $\Delta^+_4 \otimes (\Delta_3 \circ \mathrm{pr}_+)$ defines a well-defined representation of $\mathrm{SO}(4)$, and hence its associated bundle is an honest vector bundle over $M$. This proves point (1). Point (2) follows from the classical representation-theoretic facts $\Delta_3 \cong \Delta^+_4 \cong(\Delta^+_4)^*$. Finally, it is clear that, for any vector bundle equipped with a connection, the identity section of its endomorphism bundle is parallel with respect to the connection induced on the endomorphism bundle. This proves point (3). 
\end{proof}

\end{document}
